\scanpage{62}
\begin{center}6 - 45 -\hfill n\textsuperscript{o} 392\end{center}

d) Déduire de c) un exemple d'anneau de valuation discrète $A$ et d'une extension séparable $L$ de $K$ tel que $B$ ne soit de la forme $A[x]$ pour aucun $x\in B$, (partir d'une extension cyclotomique de $Q$, en utilisant l'exerc. 11 du chap. VI, § 8).

8)\quad a) Soient $A$ un anneau de Dedekind, $K$ son corps des fractions, $L$ une extension finie séparable de $K$, $B$ la fermeture intégrale de $A$ dans $L$. On suppose que pour tout idéal premier $\underline P$ de $B$, le corps résiduel de $B$ en $\underline P$ soit séparable sur celui de $A$ en $\underline P\cap A$. Montrer qu'il existe un élément $x\in B$ tel que le $B$-module $\Omega_A(B)$ des différentielles de $B$ sur $A$ soit engendré par $\mathrm{d}x$.

b) Donner un exemple d'un anneau de valuation discrète complet $A$, et d'une extension séparable $L$ du corps des fractions $K$ de $L$, tels que si $B$ est la fermeture intégrale de $A$ dans $L$, $\bar K$ et $\bar L$ les corps résiduels de $A$ et $B$ respectivement, $\bar L$ soit une extension radicielle non monogène de $\bar K$, (procéder comme dans le chap. VI, § 8, exerc. 10). Pour une telle extension, montrer que le $B$-module $\Omega_A(B)$ n'est pas monogène.

\begin{center}-:-:-:-\end{center}
\margin{61}
